% =====================================================================
%  3. Metodologia — EM ANDAMENTO
%  PRISMA 2020: critérios de elegibilidade (item 5) e seleção (item 8).
%  PENDENTE: documentar bases, datas, estratégias de busca e os detalhes
%  operacionais da triagem, incluindo revisores e resolução de divergências.
% =====================================================================
\section{Metodologia}\label{sec:metodologia}

O checklist PRISMA 2020 é utilizado como diretriz de relato da revisão
\cite{page2021prisma}. Os critérios de elegibilidade, o registro das
decisões de seleção e o uso de fontes complementares são descritos a seguir.

\subsection{Critérios de elegibilidade dos estudos}\label{sec:criterios}

Foram estabelecidos critérios de inclusão e exclusão com o objetivo de
selecionar publicações pertinentes à investigação da observabilidade
aplicada à cultura DevOps no contexto da Administração Pública.

Como critérios de inclusão, consideraram-se artigos científicos publicados
entre 2020 e 2026, nos idiomas português ou inglês, que abordassem
fundamentos, práticas, capacidades ou desafios relacionados a DevOps,
observabilidade ou transformação digital no setor público, apresentando
contribuições relevantes para os objetivos da pesquisa
(Tabela~\ref{tab:inclusao}). Os artigos devem atender a todos os critérios
pertinentes.

\begin{table}[htbp]
\centering
\caption{Critérios de inclusão}\label{tab:inclusao}
\footnotesize
\begin{tabularx}{\linewidth}{@{}lX@{}}
\toprule
\textbf{Código} & \textbf{Critério} \\
\midrule
CI1 & Artigos científicos publicados entre 2020 e 2026. \\
CI2 & Publicações em português ou inglês, com texto integral acessível. \\
CI3 & Estudos relacionados a DevOps, observabilidade ou transformação digital
      na Administração Pública. \\
CI4 & Estudos com contribuições relevantes e complementares aos eixos da
      revisão. \\
\bottomrule
\end{tabularx}
\end{table}

Foram definidos como critérios de exclusão estudos sem aderência aos eixos
temáticos, publicações restritas a soluções técnicas específicas sem
contribuição identificável para a síntese proposta, trabalhos fora do
recorte estabelecido e registros duplicados (Tabela~\ref{tab:exclusao}).

\begin{table}[htbp]
\centering
\caption{Critérios de exclusão}\label{tab:exclusao}
\footnotesize
\begin{tabularx}{\linewidth}{@{}lX@{}}
\toprule
\textbf{Código} & \textbf{Critério} \\
\midrule
CE1 & Estudos sem aderência suficiente ao problema de pesquisa. \\
CE2 & Estudos com abordagem técnica muito específica, fora do recorte
      adotado. \\
CE3 & Estudos com contribuições predominantemente sobrepostas às de outros
      trabalhos priorizados, sem acréscimo considerado essencial à síntese. \\
CE4 & Publicações fora dos critérios de período, idioma ou tipo de
      documento. \\
\bottomrule
\end{tabularx}
\end{table}

Os estudos aprovados apoiam os três eixos da fundamentação teórica:
cultura DevOps (Seção~\ref{sec:devops}), observabilidade de aplicações e
infraestrutura (Seção~\ref{sec:observabilidade}) e DevOps e observabilidade
na Administração Pública (Seção~\ref{sec:setor-publico}).

\subsection{Processo de seleção}\label{sec:processo-selecao}

A seleção foi orientada pela pertinência dos estudos à questão de pesquisa,
considerando suas contribuições conceituais, técnicas e organizacionais. As decisões documentadas
e suas justificativas são apresentadas na Seção~\ref{sec:avaliacao}.
Nesse registro, as publicações com texto integral não recuperado são
distinguidas dos estudos excluídos por inadequação ao recorte adotado.

% PRISMA 2020, item 8: o relato acima é parcial. Ainda é necessário informar
% como os critérios foram aplicados, quantos revisores participaram, se
% trabalharam independentemente e como eventuais divergências foram resolvidas.

\subsection{Uso de referências complementares}\label{sec:fontes-complementares}

Gomes, Rego e Trinta~\citeano{gomes2025} foram consultados exclusivamente
pelo resumo editorial, como referência conceitual complementar à
fundamentação teórica da Seção~\ref{sec:setor-publico}. O texto integral
não foi recuperado, de modo que o estudo permanece fora do conjunto dos
nove estudos aprovados e não é contabilizado como incluído na revisão.
As atribuições a essa fonte restringem-se às informações verificáveis no
resumo, sem análise de métodos ou resultados que dependam da leitura
integral.

\FloatBarrier
